\documentclass[11pt]{article}

\begin{document}
	
	\tableofcontents
	\newpage
	
	\setcounter{section}{2}
	\section{}
	
	\subsection{}
	\subsection*{Exercise}
	Devise three example tasks of your own that fit into the MDP framework,
	identifying for each its states, actions, and rewards. Make the three examples as different
	from each other as possible. The framework is abstract and flexible and can be applied in
	many different ways. Stretch its limits in some way in at least one of your examples.
	\subsection*{Solution}
	\begin{enumerate}
		\item \textbf{Chess Engine}: An agent which plays game of chess. States include piece positions, castling and en-passant rights. Actions are legal moves from the current position. Rewards are +1 for win, -1 for loss and 0 for draw.
		\item \textbf{Self-Driving Car}: An autonomous vehicle agent. States include car location, street position, and pedestrian presence. Actions include acceleration, braking and steering. Rewards are positive for approaching destination, negative for time passing, collisions, pedestrian hits, gas depletion, or rule violation.
		\item \textbf{Pianist}: An agent that plays the piano like a musician. States are current note positions in the musical sequence. Actions are piano key presses. Rewards measure how well the music is played.
	\end{enumerate}
	
	\subsection{}
	\subsection*{Exercise}
	Is the MDP framework adequate to usefully represent all goal-directed learning tasks? Can you think of any clear exceptions?
	\subsection*{Solution}
	No, it's not adequate to usefully represent all goal-directed learning tasks. Tasks where rewards cannot be properly defined. For example, an agent attempting to solve unsolved mathematics problems.
\end{document}